% ID: SEL-D0061-OA-E01
% ESTADO: cotejada. Segunda pasada visual y borrador compilado revisados.
% Fuente: Selectividad/sel_2010_matematicas_aplicadas/mat_apli/mat_apli_junio/mat_apli_junio/EXAMEN-ANDALUCÍA-5.pdf
% SHA-256: 5ea9fab18d8185843726e08994f39cf34f9b7dab32d86c310902a2f203dace2a
% Página: 1; opción A; ejercicio original 1.

Sea el recinto definido por las inecuaciones siguientes:
\[
x+y\leq15;\quad x\leq2y;\quad0\leq y\leq6;\quad x\geq0
\]
\begin{enumerate}[label=\alph*)]
\item \textbf{(1 punto)} Represente gráficamente dicho recinto.
\item \textbf{(1 punto)} Calcule sus vértices.
\item \textbf{(0.5 puntos)} Determine el máximo valor de la función $F(x,y)=8x+5y$ en el recinto anterior y dónde se alcanza.
\end{enumerate}
